% STAGED generated ten-item appendix. No owner approval or frozen-gate migration.
% packet-sha256: bb10ab41c5540861473f3cedf9788988e08bf1910a51c5b1608d4f092cec3f18
\section*{STAGED certificate: mobile-alpha3-r2-p2-h1}
This is a parameterized certificate display for a changed source, not a certificate for the frozen M word.
Source structure: \{"alpha":3,"h":1,"p":2,"r":2\}.

\noindent Exact semantic binding: \path{GQ2.Dyadic.GeneralTemplate.mobileEnvironment}.

\subsection*{1. Exact syntax tree and binding}
The following word is rendered from the exact Lean-exported tree. All continuous parameters remain symbolic. The named environment binds them to the actual source constructor; arbitrary assignments to these symbols do not automatically satisfy a source certificate.

% STAGED: not an approved S1.1 migration or acceptance appendix.
% Generated only from the Lean-denoting staged expression tree.
% This hash identifies a SYMBOLIC TEMPLATE, not an instantiated profinite word.
% It does NOT hash q, nuP, or the actual eta/k/correction parameter values.
% Declaration names are references, not content addresses; see the provenance sidecar.
% staged-template-hash: afa728332d136120d9856b79a03dd63b9be5f8436ac21f7e5fef185b50d1bfe3
% schema: gq2.w361.general-template.v1.staged
% source bridge: GQ2.Dyadic.GeneralTemplate.denote_mobileSplice
% exact binding: GQ2.Dyadic.GeneralTemplate.mobileEnvironment
% structural instance: alpha=3, h=1, p=2, r=2
% k_i and eta are exact Z_2 values; gamma_correction is the chosen Zhat exponent.
% The Lean bridge binds them exactly; arbitrary assignments need not satisfy source certificates.
% W_i denote exact subtrees, shared or named to bound display complexity; no simplification.
% Expanding the strictly backward references reproduces the hashed tree exactly.
% Every definition/final expression has at most 11 visible AST/ref nodes.
% Literal integer/parameter widths are not truncated; font/page layout is checked separately.
\begin{align*}
W_{1} &:= \bigl(\bigl(x_{0}\,\tau\bigr)\bigr)^{\omega_2} \\
W_{2} &:= \bigl(W_{1}\,\bigl(\bigl(\tau\bigr)^{\omega_2}\bigr)^{-1}\bigr) \\
W_{3} &:= \bigl(W_{2}\,\bigl(\sigma\bigr)^{(\mathrm{k}_{0})\omega_2}\bigr) \\
W_{4} &:= \bigl(W_{3}\bigr)^{-1} \\
W_{5} &:= \bigl(\bigl(x_{2}\,\tau\bigr)\bigr)^{\omega_2} \\
W_{6} &:= \bigl(W_{5}\,\bigl(\bigl(\tau\bigr)^{\omega_2}\bigr)^{-1}\bigr) \\
W_{7} &:= \bigl(W_{6}\,\bigl(\sigma\bigr)^{(\mathrm{k}_{2})\omega_2}\bigr) \\
W_{8} &:= \bigl(\bigl(\bigl(\sigma\bigr)^{\omega_2}\bigr)^{4}\,1\bigr) \\
W_{9} &:= \bigl(W_{7}\,W_{8}\bigr) \\
W_{10} &:= \bigl(W_{9}\bigr)^{-4}
\end{align*}
\begin{align*}
W_{11} &:= \bigl(W_{10}\,1\bigr) \\
W_{12} &:= \bigl(W_{4}\,W_{11}\bigr) \\
W_{13} &:= \bigl(\bigl(x_{1}\,\tau\bigr)\bigr)^{\omega_2} \\
W_{14} &:= \bigl(W_{13}\,\bigl(\bigl(\tau\bigr)^{\omega_2}\bigr)^{-1}\bigr) \\
W_{15} &:= \bigl(W_{14}\,\bigl(\sigma\bigr)^{(\mathrm{k}_{1})\omega_2}\bigr) \\
W_{16} &:= \bigl(\bigl(\bigl(\sigma\bigr)^{\omega_2}\bigr)^{2}\,1\bigr) \\
W_{17} &:= \bigl(W_{15}\,\bigl(1\,1\bigr)\bigr) \\
W_{18} &:= \bigl(W_{17}\bigr)^{\omega_2} \\
W_{19} &:= \bigl(W_{15}\bigr)^{-1} \\
W_{20} &:= \bigl(W_{19}\,1\bigr)
\end{align*}
\begin{align*}
W_{21} &:= \bigl(W_{18}\,W_{20}\bigr) \\
W_{22} &:= \bigl(W_{21}\bigr)^{\bigl(\bigl(\sigma\bigr)^{\omega_2}\bigr)^{16}} \\
W_{23} &:= \bigl(W_{3}\,\bigl(1\,1\bigr)\bigr) \\
W_{24} &:= \bigl(W_{23}\bigr)^{\omega_2} \\
W_{25} &:= \bigl(W_{4}\,1\bigr) \\
W_{26} &:= \bigl(W_{24}\,W_{25}\bigr) \\
W_{27} &:= \bigl(W_{26}\,1\bigr) \\
W_{28} &:= \bigl(W_{21}\,W_{27}\bigr) \\
W_{29} &:= \bigl(W_{22}\,W_{28}\bigr) \\
W_{30} &:= \bigl(W_{29}\bigr)^{\bigl(\bigl(\sigma\bigr)^{\omega_2}\bigr)^{18}}
\end{align*}
\begin{align*}
W_{31} &:= \bigl(W_{30}\,W_{27}\bigr) \\
W_{32} &:= \bigl(W_{7}\,\bigl(1\,1\bigr)\bigr) \\
W_{33} &:= \bigl(W_{32}\bigr)^{\omega_2} \\
W_{34} &:= \bigl(W_{7}\bigr)^{-1} \\
W_{35} &:= \bigl(W_{34}\,1\bigr) \\
W_{36} &:= \bigl(W_{33}\,W_{35}\bigr) \\
W_{37} &:= \bigl(W_{36}\bigr)^{\bigl(\bigl(\sigma\bigr)^{\omega_2}\bigr)^{2}} \\
W_{38} &:= \bigl(W_{37}\,1\bigr) \\
W_{39} &:= \bigl(W_{36}\,W_{38}\bigr) \\
W_{40} &:= \bigl(W_{39}\bigr)^{\bigl(\bigl(\bigl(\sigma\bigr)^{\omega_2}\bigr)^{4}\bigr)^{1}}
\end{align*}
\begin{align*}
W_{41} &:= \bigl(W_{39}\bigr)^{\bigl(\bigl(\bigl(\sigma\bigr)^{\omega_2}\bigr)^{4}\bigr)^{2}} \\
W_{42} &:= \bigl(W_{39}\bigr)^{\bigl(\bigl(\bigl(\sigma\bigr)^{\omega_2}\bigr)^{4}\bigr)^{3}} \\
W_{43} &:= \bigl(W_{39}\bigr)^{\bigl(\bigl(\bigl(\sigma\bigr)^{\omega_2}\bigr)^{4}\bigr)^{4}} \\
W_{44} &:= \bigl(W_{41}\,\bigl(W_{42}\,\bigl(W_{43}\,1\bigr)\bigr)\bigr) \\
W_{45} &:= \bigl(\bigl(W_{40}\,W_{44}\bigr)\bigr)^{\bigl(\bigl(\sigma\bigr)^{\omega_2}\bigr)^{16}} \\
W_{46} &:= \bigl(\bigl(W_{36}\bigr)^{\bigl(\bigl(\sigma\bigr)^{\omega_2}\bigr)^{4}}\,\bigl(W_{45}\,1\bigr)\bigr) \\
W_{47} &:= \bigl(W_{26}\bigr)^{-1} \\
W_{48} &:= \bigl(\bigl(\bigl(\sigma\bigr)^{\omega_2}\bigr)^{4}\bigr)^{-4} \\
W_{49} &:= \bigl(W_{48}\,1\bigr) \\
W_{50} &:= \bigl(W_{47}\,W_{49}\bigr)
\end{align*}
\begin{align*}
W_{51} &:= \bigl[\bigl(\bigl(\sigma\bigr)^{\omega_2}\bigr)^{4},\bigl(\sigma\bigr)^{\widehat{\eta}}\bigr] \\
W_{52} &:= \bigl(\bigl(x_{3}\,\tau\bigr)\bigr)^{\omega_2} \\
W_{53} &:= \bigl(W_{52}\,\bigl(\bigl(\tau\bigr)^{\omega_2}\bigr)^{-1}\bigr) \\
W_{54} &:= \bigl(\bigl(x_{4}\,\tau\bigr)\bigr)^{\omega_2} \\
W_{55} &:= \bigl(W_{54}\,\bigl(\bigl(\tau\bigr)^{\omega_2}\bigr)^{-1}\bigr) \\
W_{56} &:= \bigl[\bigl(W_{53}\,\bigl(\sigma\bigr)^{(\mathrm{k}_{3})\omega_2}\bigr),\bigl(W_{55}\,\bigl(\sigma\bigr)^{(\mathrm{k}_{4})\omega_2}\bigr)\bigr] \\
W_{57} &:= \bigl(\bigl[W_{26},W_{21}\bigr]\,\bigl(\bigl(W_{56}\,1\bigr)\,1\bigr)\bigr) \\
W_{58} &:= \bigl(W_{31}\,\bigl(\bigl(W_{26}\bigr)^{2}\,W_{57}\bigr)\bigr) \\
W_{59} &:= \bigl(\bigl(\bigl(\bigl(\sigma\bigr)^{\omega_2}\bigr)^{4}\bigr)^{8}\,\bigl(W_{51}\,W_{58}\bigr)\bigr) \\
W_{60} &:= \bigl(\bigl[W_{50},\bigl(W_{21}\,W_{16}\bigr)\bigr]\,W_{59}\bigr)
\end{align*}
\begin{align*}
W_{61} &:= \bigl(W_{46}\,\bigl(\bigl(W_{50}\bigr)^{2}\,W_{60}\bigr)\bigr) \\
W_{62} &:= \bigl(\bigl[W_{9},\bigl(\sigma\bigr)^{\widehat{\eta}}\bigr]\,\bigl(W_{31}\,W_{61}\bigr)\bigr) \\
W_{63} &:= \bigl(\bigl[W_{12},\bigl(W_{15}\,W_{16}\bigr)\bigr]\,\bigl(\bigl(W_{9}\bigr)^{8}\,W_{62}\bigr)\bigr) \\
W_{64} &:= \bigl(\bigl(\bigl(W_{12}\bigr)^{2}\,W_{63}\bigr)\,\bigl(\bigl(1\bigr)^{\gamma_{\mathrm{correction}}}\bigr)^{-1}\bigr) \\
W_{65} &:= \bigl(\bigl(\bigl(\sigma\bigr)^{\omega_2}\bigr)^{-4}\,1\bigr) \\
W_{66} &:= \bigl(\bigl(x_{0}\bigr)^{-1}\,W_{65}\bigr) \\
W_{67} &:= \bigl(\bigl(x_{2}\,\bigl(\tau\,1\bigr)\bigr)\bigr)^{\omega_2} \\
W_{68} &:= \bigl(\bigl(\bigl(x_{2}\bigr)^{\sigma}\bigr)^{-1}\,\bigl(W_{67}\,1\bigr)\bigr) \\
W_{69} &:= \bigl(x_{1}\,\bigl(\tau\,1\bigr)\bigr) \\
W_{70} &:= \bigl(W_{69}\bigr)^{\omega_2}
\end{align*}
\begin{align*}
W_{71} &:= \bigl(\bigl(x_{1}\bigr)^{-1}\,1\bigr) \\
W_{72} &:= \bigl(W_{70}\,W_{71}\bigr) \\
W_{73} &:= \bigl(x_{0}\,\bigl(\tau\,1\bigr)\bigr) \\
W_{74} &:= \bigl(W_{73}\bigr)^{\omega_2} \\
W_{75} &:= \bigl(\bigl(x_{0}\bigr)^{-1}\,1\bigr) \\
W_{76} &:= \bigl(W_{74}\,W_{75}\bigr) \\
W_{77} &:= \bigl(\bigl(W_{76}\bigr)^{\bigl(\bigl(\sigma\bigr)^{\omega_2}\bigr)^{4}}\,\bigl(W_{76}\,1\bigr)\bigr) \\
W_{78} &:= \bigl(\bigl(W_{72}\bigr)^{\bigl(\bigl(\sigma\bigr)^{\omega_2}\bigr)^{4}}\,W_{77}\bigr) \\
W_{79} &:= \bigl(\bigl(W_{72}\bigr)^{\bigl(\bigl(\sigma\bigr)^{\omega_2}\bigr)^{8}}\,W_{78}\bigr) \\
W_{80} &:= \bigl(\bigl(\bigl[x_{3},x_{4}\bigr]\,1\bigr)\,1\bigr)
\end{align*}
\begin{align*}
W_{81} &:= \bigl(\bigl(\bigl(\sigma\bigr)^{\omega_2}\bigr)^{8}\,\bigl(W_{68}\,\bigl(W_{79}\,W_{80}\bigr)\bigr)\bigr) \\
W_{82} &:= \bigl(\bigl(W_{66}\bigr)^{2}\,\bigl(\bigl[W_{66},x_{1}\bigr]\,W_{81}\bigr)\bigr) \\
W_{83} &:= \bigl(W_{2}\bigr)^{-1} \\
W_{84} &:= \bigl(W_{83}\,W_{65}\bigr) \\
W_{85} &:= \bigl(\bigl(W_{6}\,\bigl(1\,1\bigr)\bigr)\bigr)^{\omega_2} \\
W_{86} &:= \bigl(\bigl(\bigl(W_{6}\bigr)^{\sigma}\bigr)^{-1}\,\bigl(W_{85}\,1\bigr)\bigr) \\
W_{87} &:= \bigl(W_{14}\,\bigl(1\,1\bigr)\bigr) \\
W_{88} &:= \bigl(W_{87}\bigr)^{\omega_2} \\
W_{89} &:= \bigl(W_{14}\bigr)^{-1} \\
W_{90} &:= \bigl(W_{89}\,1\bigr)
\end{align*}
\begin{align*}
W_{91} &:= \bigl(W_{88}\,W_{90}\bigr) \\
W_{92} &:= \bigl(W_{2}\,\bigl(1\,1\bigr)\bigr) \\
W_{93} &:= \bigl(W_{92}\bigr)^{\omega_2} \\
W_{94} &:= \bigl(W_{83}\,1\bigr) \\
W_{95} &:= \bigl(W_{93}\,W_{94}\bigr) \\
W_{96} &:= \bigl(\bigl(W_{95}\bigr)^{\bigl(\bigl(\sigma\bigr)^{\omega_2}\bigr)^{4}}\,\bigl(W_{95}\,1\bigr)\bigr) \\
W_{97} &:= \bigl(\bigl(W_{91}\bigr)^{\bigl(\bigl(\sigma\bigr)^{\omega_2}\bigr)^{4}}\,W_{96}\bigr) \\
W_{98} &:= \bigl(\bigl(W_{91}\bigr)^{\bigl(\bigl(\sigma\bigr)^{\omega_2}\bigr)^{8}}\,W_{97}\bigr) \\
W_{99} &:= \bigl(\bigl(\bigl[W_{53},W_{55}\bigr]\,1\bigr)\,1\bigr) \\
W_{100} &:= \bigl(\bigl(\bigl(\sigma\bigr)^{\omega_2}\bigr)^{8}\,\bigl(W_{86}\,\bigl(W_{98}\,W_{99}\bigr)\bigr)\bigr)
\end{align*}
\begin{align*}
W_{101} &:= \bigl(\bigl(W_{84}\bigr)^{2}\,\bigl(\bigl[W_{84},W_{14}\bigr]\,W_{100}\bigr)\bigr) \\
W_{102} &:= \bigl(\bigl(W_{64}\,W_{82}\bigr)\,\bigl(W_{101}\bigr)^{-1}\bigr) \\
R_{\mathrm{mobile}} &:= W_{102}
\end{align*}

\paragraph{Checked declarations}
\begin{itemize}
\item \path{GQ2.Dyadic.GeneralTemplate.denote_mobileSplice}
\item \path{GQ2.Dyadic.GeneralTemplate.Tree.eval_eq_pinned}
\end{itemize}

\subsection*{2. Tame and pro-two specialization reductions}
Write W=Shadow(Rcorr) H Shadow(H)\textasciicircum{}(-1), with the literal corrected normalized word Rcorr and compact template H from item 1. On pro-two targets with tau=1, shadow is the identity, the prime-to-two correction vanishes, and the proved Nielsen normalization transports the complete presented core. On the canonical tame marking, the correction removes the actual defect and the splice specializes to one.

The mobile core certificate retains the exact actual-unit equation for nuP(D), with arbitrary actual eta. When obtained from field selection, eta also retains that selector's proved congruence; it is not replaced by the representative.

\paragraph{Checked declarations}
\begin{itemize}
\item \path{GQ2.Dyadic.eval_pro2_tameSplice}
\item \path{GQ2.Dyadic.tameSpecializes_tameSplice}
\item \path{GQ2.Dyadic.Instances.MpcMobileSpliceSupply.mpcMobileTameSplice_core}
\item \path{GQ2.Dyadic.Instances.MpcMobileNuSource.mpcMobileTameCorrectedW_specializes}
\end{itemize}

\subsection*{3. Evaluated Fox Jacobian}
All three profiles use a finite base group C and finite exponent-two coefficient module A. These are identities on all cochains, not just cycles. Columns not displayed vanish. The resolver is pinned at 4 times the finite base-group exponent and is valid at WordLift level.

Scalar: every generator action is trivial. Unramified: finite exponent-two coefficients, wild and tau actions trivial, sigma fixed-point-free with trivial two-primary action, and death of the actual pinned changed-word family. Ramified: explicit finite surjective head projection with compatible action, wild head one, actual tame relation, and tau fixed-point-free with trivial two-primary action. No raw-word death or commuting sigma/tau premise is inserted.

\paragraph{scalar}
\[\begin{pmatrix}x_{\tau}\\x_{\tau}\end{pmatrix}\]
\paragraph{unramified}
\[\begin{pmatrix}S^{-1}\bigl(x_{\tau}\bigr)\\\bigl(x_{\tau}+\bigl(I+\bigl(-S^{-1}\bigr)\bigr)\bigl(x_{2}\bigr)\bigr)\end{pmatrix}\]
\paragraph{ramified}
\[\begin{pmatrix}\bigl(\bigl(S^{-1}\circ \bigl(T+\bigl(-I\bigr)\bigr)\bigr)\bigl(x_{\sigma}\bigr)+\bigl(S^{-1}+\bigl(-\left(\sum_{j=0}^{q-1}T^{j}\right)\bigr)\bigr)\bigl(x_{\tau}\bigr)\bigr)\\\bigl(-S^{-1}\bigr)\bigl(x_{2}\bigr)\end{pmatrix}\]
\paragraph{Checked declarations}
\begin{itemize}
\item \path{GQ2.Dyadic.Instances.TameSpliceScalarFoxNormal.mobileSplice_scalar_rows}
\item \path{GQ2.Dyadic.Instances.TameSpliceUnramifiedFoxRows.mobileSplice_unramified_rows}
\item \path{GQ2.Dyadic.Instances.TameSpliceRamifiedFoxNormal.tameSplice_ramified_rows}
\item \path{GQ2.Dyadic.GeneralFoxTemplate.scalar_mobile_input}
\item \path{GQ2.Dyadic.GeneralFoxTemplate.unramified_mobile_input}
\item \path{GQ2.Dyadic.GeneralFoxTemplate.ramified_mobile_input}
\end{itemize}

\subsection*{4. Ordered invertible operations and carried inverses}
The following lists are exported from the same typed profile as the input and target rows. Their denotations are definitionally the lists in the original FoxCertificate records. Every scale carries a formal inverse; the actual records prove both inverse products.

In each nonscalar sector k is the half-period of the actual coefficient action, so its odd order is 2k+1. It is not the handle rank, not an assumed ambient group order, and not a fixed source parameter. q is the source residue size. Composition applies the inner operator first.

\paragraph{scalar}
\[\begin{aligned}\mathcal R_{1}&:\quad R_2 \leftarrow R_2+\bigl(-I\bigr)\circ R_1\\\mathcal C&=\mathrm{id}\end{aligned}\]
\paragraph{unramified}
\[\begin{aligned}\mathcal R_{1}&:\quad R_1 \leftarrow S\circ R_1\quad[\text{carried inverse: }S^{-1}]\\\mathcal R_{2}&:\quad R_2 \leftarrow R_2+\bigl(-I\bigr)\circ R_1\\\mathcal R_{3}&:\quad R_2 \leftarrow \left(\sum_{j=0}^{k-1}S^{-(2j+1)}\right)\circ R_2\quad[\text{carried inverse: }\bigl(I+\bigl(-S^{-1}\bigr)\bigr)]\\\mathcal C&=\mathrm{id}\end{aligned}\]
\paragraph{ramified}
\[\begin{aligned}\mathcal R_{1}&:\quad R_1 \leftarrow \bigl(\left(\sum_{j=0}^{k-1}T^{2j+1}\right)\circ S\bigr)\circ R_1\quad[\text{carried inverse: }\bigl(S^{-1}\circ \bigl(T+\bigl(-I\bigr)\bigr)\bigr)]\\\mathcal R_{2}&:\quad R_2 \leftarrow \bigl(-S\bigr)\circ R_2\quad[\text{carried inverse: }\bigl(-S^{-1}\bigr)]\\\mathcal C_{1}&:\quad x_{\sigma}\leftarrow x_{\sigma}+\bigl(-\bigl(\bigl(\left(\sum_{j=0}^{k-1}T^{2j+1}\right)\circ S\bigr)\circ \bigl(S^{-1}+\bigl(-\left(\sum_{j=0}^{q-1}T^{j}\right)\bigr)\bigr)\bigr)\bigr)\bigl(x_{\tau}\bigr)\end{aligned}\]
\paragraph{Checked declarations}
\begin{itemize}
\item \path{GQ2.Dyadic.GeneralFoxTemplate.scalar_mobile_signature}
\item \path{GQ2.Dyadic.GeneralFoxTemplate.unramified_mobile_signature}
\item \path{GQ2.Dyadic.GeneralFoxTemplate.ramified_mobile_signature}
\end{itemize}

\subsection*{5. Reduced normal matrices}
Replay of the preceding invertible lists gives these literal identity-pivot targets on all cochains. The scalar second row is zero. The ramified sigma pivot is obtained after the displayed column shear; it must not be read as an identity for the unsheared input.

\paragraph{scalar}
\[\begin{pmatrix}x_{\tau}\\0\end{pmatrix}\]
\paragraph{unramified}
\[\begin{pmatrix}x_{\tau}\\x_{2}\end{pmatrix}\]
\paragraph{ramified}
\[\begin{pmatrix}x_{\sigma}\\x_{2}\end{pmatrix}\]
\paragraph{Checked declarations}
\begin{itemize}
\item \path{GQ2.Dyadic.GeneralFoxTemplate.scalar_mobile_signature}
\item \path{GQ2.Dyadic.GeneralFoxTemplate.unramified_mobile_signature}
\item \path{GQ2.Dyadic.GeneralFoxTemplate.ramified_mobile_signature}
\end{itemize}

\subsection*{6. Stokes chain identity and all coefficients}
The actual word supplies a single coefficient-independent UniformStokesSupply. It includes pinned resolution, endpoint and full Stokes duality for every finite elementary coefficient module pushed from the actual source. The simple-action profiles feed the proved composition-series exact-diagram extension; checking dimensions only on simple modules is not substituted for that argument. The statement is uniform in all source parameters and h.

\paragraph{Checked declarations}
\begin{itemize}
\item \path{GQ2.Dyadic.Instances.MpcMobileSpliceSupply.mpcMobileTameSplice_uniformStokesSupply}
\end{itemize}

\subsection*{7. Full extraspecial Hessian}
The exported quadratic expression below is the full endpoint phase. Here q\_0 is the coefficient quadratic form and B=polar(q\_0), distinct from the residue size q in the Fox rows. It is not inferred from its polar matrix. The unramified comparison retains the factor-set charge until its proved faithful-head vanishing; the ramified splice retains the complete extraspecial evaluation and then uses the existing head swap. This is not a claim that the raw ramified expression equals the displayed form before changing coordinates.

\[Q=\bigl(\bigl(q_0(c_0)+B(c_0,c_1)\bigr)+\bigl(B(u_{0},v_{0})+0\bigr)\bigr).\]
\paragraph{Checked declarations}
\begin{itemize}
\item \path{GQ2.Dyadic.hessEvalZ_tameSplice_ramified}
\item \path{GQ2.Dyadic.GeneralPhaseTemplate.denote_compactPhase}
\item \path{GQ2.Dyadic.GeneralPhaseTemplate.sourcePhase_eq_denote}
\item \path{GQ2.Dyadic.Instances.MpcMobileUnramifiedPhase.hessRelZ_mpcUnit_unramified_compactPhase_of_sigma2_one}
\end{itemize}

\subsection*{8. Linear normal coordinates and polar inverse}
The original source comparison map is proved additive, then upgraded to an F2-linear equivalence with exactly the same forward and inverse functions. Its domain is the actual source cocycle quotient. The original Hessian record proves QZeroBar(x)=Q(cov(x)), and the two inverse equations. The matrices below are relative to the coefficient pairing B; B itself is not declared to be the identity matrix.

\[J=\begin{pmatrix}1&1&0&0\\1&0&0&0\\0&0&0&1\\0&0&1&0\end{pmatrix},\qquad J^{-1}=\begin{pmatrix}0&1&0&0\\1&1&0&0\\0&0&0&1\\0&0&1&0\end{pmatrix}.\]
Coordinate order: $c_0$, $c_1$, $u_{0}$, $v_{0}$.
\paragraph{Checked declarations}
\begin{itemize}
\item \path{GQ2.Dyadic.Instances.WordPhaseCompactCertificates.sourceCompactHessianOfComparison_cov_apply}
\item \path{GQ2.Dyadic.GeneralPhaseTemplate.sourceHessian_cov_eq_denote}
\item \path{GQ2.Dyadic.GeneralPhaseTemplate.polar_denote_compactPhase}
\item \path{GQ2.Dyadic.Instances.WordPhasePolarNormal.reducedCoefficientMatrix_mul_inverse}
\item \path{GQ2.Dyadic.Instances.TameSpliceBoundaryHessianCertificates.mobileBoundaryHessian}
\end{itemize}

\subsection*{9. Original affine phase-cover record}
The original Hessian record contains the original PhaseCoverCertificate. For |V|=2\textasciicircum{}d, its base dimension is (h+1)d and its base sign is +1. The diagonal and cross embeddings put v into the first and second head slots respectively and zero in every other slot. It proves the full diagonal identity Q(j0(v))=dat.f(v,v), the cross-polar identity, the endpoint cardinality and the signed Gauss sum.

At every relevant actual boundary lift, the named boundary producer derives the projected head hypotheses internally and returns this original record. The downstream affine determinant remains a theorem for every required cover and character, not just an unshifted Gauss value.

% STAGED symbolic phase-cover data; not a source-word or coefficient-space hash.
% cover-template-hash: d21b7952f6b4eeccbf62573e6bd1b5f258e999854fe27041ccd02dbc5b2597f9
% phase-template-hash: e48eecf746e3bb42d6f28431b8d8ecfe5839d018b4afdd063e222254e28090cc
% exact bridge: GQ2.Dyadic.GeneralPhaseCoverTemplate.compactTemplate_agrees
% q_0 is the coefficient quadratic form, not residue cardinality q; #V=2^d.
% Embedding vectors are in the listed flattened head/handle coordinate order.
\[j_0(v)=\bigl(v,0,0,0\bigr),\qquad j_1(v)=\bigl(0,v,0,0\bigr).\]
\[\operatorname{diag}(v)=q_0(v),\qquad m=\bigl(\bigl(1+1\bigr)\cdot d\bigr),\qquad \varepsilon=1.\]
\[Q=\bigl(\bigl(q_0(c_0)+B(c_0,c_1)\bigr)+\bigl(B(u_{0},v_{0})+0\bigr)\bigr).\]
% Under compactTemplate_agrees hypotheses, diag(v)=dat.f(v,v);
% the original cover certifies #W=2^(2m), Gauss(Q)=epsilon*2^m and the polar/kappa identities.
% coordinate order: [{"coordinate":"head0"},{"coordinate":"head1"},{"coordinate":"handleU","index":0},{"coordinate":"handleV","index":0}]

\[Q(j_0(v))=\operatorname{dat}.f(v,v),\qquad B_Q(j_0(v),j_1(w))=\operatorname{dat}.f(v,w)+\operatorname{dat}.f(w,v).\]
\paragraph{Checked declarations}
\begin{itemize}
\item \path{GQ2.Dyadic.Instances.WordPhaseCompactCertificates.compactPhaseCover}
\item \path{GQ2.Dyadic.GeneralPhaseTemplate.compactPhaseCover_tree}
\item \path{GQ2.Dyadic.GeneralPhaseCoverTemplate.compactTemplate_agrees}
\item \path{GQ2.Dyadic.Instances.TameSpliceBoundaryHessianCertificates.mobileBoundaryHessian}
\item \path{GQ2.Dyadic.Instances.MpcMobileSpliceBoundary.mpcMobileTameSplice_determinant}
\end{itemize}

\subsection*{10. Same-tree hashes, replay and trust boundary}
The word, Fox, phase and phase-cover hashes identify symbolic templates. They do not hash actual q, nuP, eta, shear, correction, coefficient-space or action-period values. Uniform Lean denotation and certificate-signature bridges bind their semantic uses. This combined packet additionally binds the generated section content and direct source bytes.

Strong checking re-emits all four Lean exports and checks every cited declaration's existence and standard-three axiom footprint, and compares the printed type transcript on replay. It does not assert independently authored expected signatures. The probe imports existing olean artifacts; direct source hashes alone do not prove that those bytes produced those artifacts or came from an atomic filesystem snapshot. A separate isolated build is required. External dependency artifacts are reused. Prose remains an editorial exposition of the checked objects. This STAGED packet does not approve the changed source selection, replace the frozen five-fixture manifest, activate a merge gate, or certify the original frozen-M theorem.

word template: \path{afa728332d136120d9856b79a03dd63b9be5f8436ac21f7e5fef185b50d1bfe3}.
Fox template: \path{b2b4dc5c45cbf45f0fb96d9cd6a325fd5a4c58781a64dfcde6b4a805a649a7e7}.
phase template: \path{e48eecf746e3bb42d6f28431b8d8ecfe5839d018b4afdd063e222254e28090cc}.
phase cover template: \path{d21b7952f6b4eeccbf62573e6bd1b5f258e999854fe27041ccd02dbc5b2597f9}.
\par Combined packet: \path{bb10ab41c5540861473f3cedf9788988e08bf1910a51c5b1608d4f092cec3f18}.
